\documentclass[11pt,a4paper]{article}

%-------------------------------------------
% Paquetes básicos
%-------------------------------------------
\usepackage[spanish]{babel}
\usepackage[utf8]{inputenc}
\usepackage[T1]{fontenc}
\usepackage{lmodern}
\usepackage{amsmath, amssymb, amsthm}
\usepackage{geometry}
\usepackage{enumitem}
\usepackage{hyperref}
\usepackage{braket}

\geometry{margin=2.5cm}

%-------------------------------------------
% Entorno para ejercicios
%-------------------------------------------
\newtheoremstyle{ejercicio}
  {10pt}   % espacio arriba
  {10pt}   % espacio abajo
  {\normalfont}
  {}
  {\bfseries}
  {.}
  {0.5em}
  {}

\theoremstyle{ejercicio}
\newtheorem{ejercicio}{Ejercicio}
\newtheorem{theorem}{Teorema}
\newtheorem{definition}{Definición}

%-------------------------------------------
\title{Guía de Ejercicio II:\\Oda a la Aproximación}
\author{nICO no IKO 777}

\date{}



% -------------------------------------------------
% DOCUMENTO
% -------------------------------------------------
\begin{document}

\maketitle
\textit{No tener moderación muchas veces es causa de que el bien se convierta en mal y la virtud en vicio.}  San Ignacio de Loyola

%\vspace{0.5cm}

\section*{Indicaciones para el desarrollo y la entrega}

\begin{itemize}

\item \textbf{Formato de entrega.}  
La entrega debe realizarse exclusivamente en formato PDF,
generado a partir de un documento escrito en \LaTeX.
No se aceptarán entregas en formato Word u otros editores.

\item \textbf{Presentación.}  
Sea ordenado y claro en la redacción.  
Incluya demostraciones completas en los ejercicios teóricos
y explique adecuadamente los resultados obtenidos en los ejercicios prácticos.
Las gráficas deben estar correctamente rotuladas y comentadas.

\item \textbf{Plazo de entrega.}  
El plazo es de \textbf{dos semanas contadas a partir del \today}.  
La fecha límite de entrega es el \textbf{17/03/2026}. 

\item \textbf{Uso de consultas.}  
Pueden realizar consultas durante el desarrollo de la guía.
Con todo gusto los orientaré.

\item \textbf{Clase de apoyo.}  
El martes 24 utilizaremos la clase para resolver ejercicios
y discutir dificultades comunes.  
Si es necesario, podremos organizar una clase adicional optativa
para profundizar en los temas más complejos.

\item \textbf{Uso de Inteligencia Artificial.}  
Se permite el uso de herramientas de IA como apoyo.
Sin embargo, deben ser plenamente conscientes de lo que están haciendo.
La IA puede cometer errores o generar argumentos incorrectos.
Usted es responsable del contenido que entrega.

\end{itemize}

\begin{enumerate}

% -------------------------------------------------
\item 
Se tienen los datos:

\[
(0,1),\quad (1,2),\quad (2,1),\quad (3,3).
\]

\begin{enumerate}
\item Construya el polinomio interpolante de Lagrange.
\item Construya el ajuste lineal por mínimos cuadrados.
\item Compare:
\begin{itemize}
\item Grado del polinomio
\item Error en los nodos
\item Comportamiento fuera del intervalo
\end{itemize}
\item ¿Cuál método es más estable ante pequeñas perturbaciones en los datos?
\item Construya el polinomio interpolante de Newton con $(1,2),\quad (2,4.1),\quad (3,5.9),\quad (4,8.2).$ {\bf datos nuevos.} Compare la regresion lineal y Newton
\end{enumerate}


% -------------------------------------------------
\item 
Se dan datos de una almacenaje

\begin{table}[h]
\centering
\begin{tabular}{c|cc}
\hline
$t$ & $entrada$ & $salida$ \\
\hline
0 & 420.157651 & 255.839425 \\
1 & 541.411707 & 310.958667 \\
2 & 901.157457 & 485.680014 \\
3 & 639.080229 & 402.997356 \\
4 & 750.875606 & 495.560775 \\
\hline
\end{tabular}
\caption{Datos experimentales}
\end{table}

\begin{enumerate}
\item Construya un spline cúbico interpolante.
\item Ajuste un modelo logístico.
\item Compare:
\begin{itemize}
\item Comportamiento en extrapolación.
\item Sensibilidad ante perturbaciones pequeñas $10^{-3}$
\item Discuta
\end{itemize}
\end{enumerate}

% -------------------------------------------------
%\item \textbf{Aproximación Racional tipo Padé vs Polinomio de Taylor}

%Sea

%\[
%f(x)=\ln(1+x).
%\]

%\begin{enumerate}
%\item Construya el polinomio de Taylor de orden 3 en $x=0$.
%\item Construya el aproximante racional de Padé $(1,1)$.
%\item Compare los errores en el intervalo $[-0.5,0.5]$.
%\item Discuta cuál aproxima mejor cerca del borde del intervalo.
%\end{enumerate}

% -------------------------------------------------
%\item \textbf{Base Ortogonal vs Base Monomial}
%
%Considere las funciones
%
%\[
%1,\quad x,\quad x^2
%\]
%
%en el intervalo $[-1,1]$.
%
%\begin{enumerate}
%\item Ortogonalice el conjunto usando el producto interno
%\[
%\int_{-1}^1 f(x)g(x)\,dx.
%\]
%\item Compare la matriz asociada con la obtenida usando la base monomial.
%\item Explique por qué la base ortogonal mejora el condicionamiento.
%\end{enumerate}



% -------------------------------------------------
\item 

%\textbf{Ortogonalidad mediante Integración por Partes}

Sea
\[
p_n(x)=\frac{1}{2^n n!}\frac{d^n}{dx^n}\left[(x^2-1)^n\right],
\]
con $x\in[-1,1]$. Los polinomios de Legendre.

Demuestre que
con
\[
\braket{p_n|p_m}= \int_{-1}^{1} p_n(x) p_m(x)\,dx 
\quad, n\neq m\]

Los polinomios de Legendre son ortogonales en $[-1,1]$. \textit{Sugerencia: Entienda el problema 
para $n=3$}.

% -------------------------------------------------




\item Considere la siguiente definición:
% =====================================================
\begin{definition}[Regresión $n-$esima (Cuadrática $n=2$)]
Sea un conjunto de datos
\[
(x_i,y_i), \quad i=1,\dots,m.
\]

\noindent La regresión de orden $n$ (cuadrática si $n=2$) consiste en aproximar los datos
mediante un modelo polinómico de grado $n$ (dos, $n=2$):

\[
y \approx {\bf c_0 + c_1 x + c_2 x^2}+ \cdots c_n x^n
\]

Los coeficientes $c_0,c_1,c_2, \dots c_n$ se determinan
minimizando el error cuadrático total

\[
\sum_{i=1}^m 
\left( y_i - y \right)^2.
\]

En forma matricial, el problema se escribe como

\[
\min_{\mathbf{c}}
\|A\mathbf{c}-\mathbf{y}\|_2^2,
\]

donde la matriz $A$ se denomina matriz de diseño.
\end{definition}

Con esto en mente, considera los siguientes datos experimentales:

\[
(0,1.2), \quad (0.5,1.8), \quad (1,2.5), \quad (1.5,2.1), \quad (2,3.2).
\]

\subsection*{Parte I}
\begin{enumerate}
\item Construya el polinomio interpolante usando Newton.
\item Construya el spline cúbico interpolante natural.
\item Compare gráficamente:
\item Analice el comportamiento fuera del intervalo $[0,2]$.
\item Ajuste un modelo de regresión cúbica mediante mínimos cuadrados.
\item Compare el error cuadrático total del modelo cúbico con el error del polinomio interpolante.
\end{enumerate}


\subsection*{Parte II}

\begin{enumerate}
\item Considere el polinomio interpolante obtenido en la Parte I.
\begin{itemize}
\item Calcule su expansión de Taylor de orden $3$ centrada en $x=1$.
\item Construya el aproximante de Padé $(1,1)$ en torno a $x=1$.
\end{itemize}
\item Repita el procedimiento anterior utilizando ahora el modelo de regresión cúbica como función base.
\item Compare los errores en el intervalo $[0.8,1.2]$ para:
\begin{itemize}
\item Taylor basado en interpolación,
\item Padé basado en interpolación,
\item Taylor basado en regresión,
\item Padé basado en regresión.
\end{itemize}
\item ¿Influye el hecho de que la función base provenga de interpolación o regresión?
\end{enumerate}



\end{enumerate}

\textbf{EXTRA, extra ... (No es necesario entregar)}

Sea

\[
p_n(x)=\frac{d^n}{dx^n}\left[(x^2-1)^n\right].
\]

Demuestre que

\[
\int_{-1}^{1} p_n(x) q(x)\,dx=0
\]

para todo polinomio $q$ de grado menor que $n$.

\textit{Sugerencia: aplicar integración por partes repetidamente.}\\


\textit{Las matemáticas parecen dotar a uno de nuevo sentido.}
Charles Darwin

\vfill
\begin{center}
\textit{La interpolación impone coincidencia exacta;  
la aproximación busca equilibrio global.}
\end{center}

\end{document}